\documentclass[11pt]{article}
\usepackage[review]{acl}
\usepackage{times}
\usepackage{latexsym}
\usepackage[T1]{fontenc}
\usepackage[utf8]{inputenc}
\usepackage{microtype}
\usepackage{inconsolata}
\usepackage{graphicx}
\usepackage{amsmath}
\usepackage{amssymb}
\usepackage{booktabs}
\usepackage{multirow}
\usepackage{xcolor}
\usepackage{subcaption}

\newcommand{\method}{SafeHead}
\newcommand{\lam}{\lambda}
\newcommand{\remark}[1]{{\color{red}\textbf{[#1]}}}
\newcommand{\better}[1]{{\color{green!60!black}\scriptsize{#1}}}
\newcommand{\worse}[1]{{\color{red!70!black}\scriptsize{#1}}}

\title{Inference-Time Safety Recovery for Fine-Tuned Language Models \\via Activation-Conditioned Hypernetworks}

\author{First Author \\
  Affiliation \\
  \texttt{email@domain} \\\And
  Second Author \\
  Affiliation \\
  \texttt{email@domain} \\}

\begin{document}
\maketitle

% ===========================================================================
% ABSTRACT
% ===========================================================================
\begin{abstract}
Safety alignment in large language models is fragile under downstream fine-tuning: even benign task adaptation can substantially degrade refusal behavior. Existing mitigation strategies require per-model retraining, access to safety data, or model-specific guardrails that do not scale to the growing ecosystem of fine-tuned checkpoints. We propose \method{}, a framework that generates a lightweight safety classifier for each fine-tuned model at inference time, without any additional training. Our key idea is to treat intermediate activations of a fine-tuned model as a compact behavioral fingerprint and train a hypernetwork that maps these activations to the parameters of a small side network acting as a prompt-level safety head. The hypernetwork employs layer-wise factorized weight generation with a direction-magnitude activation decomposition, producing model-specific safety heads from only $N_\text{cal}$ calibration prompts. We evaluate on Qwen2-7B and LLaMA-3-8B across 22 fine-tuning domains with three multi-holdout splits. Generated safety heads reduce harmful output rates from up to 43\% to near 0\% on unseen domains while preserving downstream task accuracy within 1\%, and transfer zero-shot to AdvBench (100\%) and HEx-PHI (97--100\%) benchmarks never seen during training.
\end{abstract}

% ===========================================================================
% 1. INTRODUCTION
% ===========================================================================
\section{Introduction}
\label{sec:introduction}

Large language models (LLMs) are typically released with safety alignment that reduces harmful outputs, achieved through instruction tuning and preference optimization \cite{ouyang2022training, bai2022training}. However, a growing body of evidence shows that this alignment is brittle: downstream fine-tuning---including seemingly benign adaptation---can significantly degrade refusal behavior and increase vulnerability to harmful prompts \cite{qi2024finetuning, qi2025tokensdeep, fraser2025finetuninglowers}. Even a small number of fine-tuning steps on harmless data can erode safety properties that required substantial effort to instill \cite{huang2024lisa}. This fragility presents a pressing challenge: the open-weight ecosystem increasingly relies on fine-tuning, adapters, and community-contributed checkpoints, yet safety guarantees do not transfer reliably across these model variants.

Existing approaches to restoring safety after fine-tuning fall into two categories. \textit{Post-hoc realignment} methods apply additional safety training to each fine-tuned checkpoint \cite{huang2024lisa, hsu2024safe}, requiring access to safety data and incurring per-model training costs. \textit{External safety classifiers} such as Llama Guard \cite{inan2023llama} or prompt shields \cite{kim2024robust} operate independently of the target model and thus cannot account for how a particular fine-tuning recipe has altered the model's safety boundaries. Activation-based detectors \cite{jiang2025hiddendetect} do leverage model internals but require fitting per-model probes that do not transfer to new checkpoints.

We start from a different observation: \textbf{prompt safety is model-dependent}. The same input prompt may be handled safely by one fine-tuned variant yet elicit harmful compliance from another, depending on the domain, data mixture, and training duration. This suggests that an effective safety mechanism must be conditioned on the specific fine-tuned model, not just on the prompt text.

Building on this insight, we propose \method{}, a framework for \textit{activation-conditioned safety head generation}. We train a hypernetwork \cite{ha2017hypernetworks} that takes intermediate activations from a fine-tuned LLM as input and outputs the weights of a lightweight side network---a prompt-level safety classifier we call a \textit{safety head}. The safety head follows the Ladder Side-Tuning architecture \cite{sung2022lst}, operating alongside the frozen backbone through gated lateral connections. At deployment, given any new fine-tuned checkpoint, we extract activations from a small set of $N_\text{cal}$ calibration prompts, run them through the hypernetwork, and obtain a model-specific safety head---all in a single forward pass, with no additional training.

A key technical challenge is that the mapping from activations to safety head weights is high-dimensional: the side network contains approximately 237M parameters for Qwen2-7B. Generating all weights from a handful of activation vectors is severely ill-conditioned. We address this through two design choices: (1) \textit{layer-wise factorization}, where each side network layer's weights are predicted independently from corresponding backbone activations, reducing the effective output dimensionality; and (2) \textit{direction-magnitude decomposition}, which splits activation vectors into unit directions (capturing \textit{what} changed during fine-tuning) and log-magnitudes (capturing \textit{how much}), processed by separate specialized encoders.

% \remark{Figure 1 here: pipeline overview}

\paragraph{Contributions.}
\begin{itemize}
    \item We formalize \textbf{model-conditioned safety assessment}: prompt safety depends on the specific fine-tuned checkpoint, and the safety classifier should be generated as a function of the model's behavioral fingerprint.
    \item We propose \textbf{\method{}}, a hypernetwork that maps fine-tuned model activations to safety head weights via layer-wise factorized generation with direction-magnitude decomposition, enabling zero-shot instantiation of safety classifiers for unseen models.
    \item We evaluate across \textbf{two model architectures} (Qwen2-7B, LLaMA-3-8B), \textbf{22 fine-tuning domains}, and \textbf{three safety benchmarks}, demonstrating that generated safety heads recover harmful detection ($>$99\% on BeaverTails, 100\% on AdvBench, 97--100\% on HEx-PHI) while preserving task performance ($<$1\% accuracy drop) on domains never seen during hypernetwork training.
\end{itemize}


% ===========================================================================
% 2. RELATED WORK
% ===========================================================================
\section{Related Work}
\label{sec:related_work}

\paragraph{Safety Degradation under Fine-Tuning.}
\citet{qi2024finetuning} demonstrate that fine-tuning aligned LLMs on benign data compromises safety, even without malicious intent. \citet{qi2025tokensdeep} show that alignment is shallow, encoded in only a few token positions. \citet{fraser2025finetuninglowers} confirm that fine-tuning lowers safety across multiple evaluation dimensions. Several defenses have been proposed: LISA \cite{huang2024lisa} applies lazy alignment during fine-tuning, Safe LoRA \cite{hsu2024safe} projects LoRA updates to preserve alignment, and Vaccine \cite{yi2024vaccine} adds perturbation-aware alignment. All require intervention during or after fine-tuning for each checkpoint individually.

\paragraph{Safety Classifiers and Guardrails.}
External classifiers like Llama Guard \cite{inan2023llama} and WildGuard \cite{han2024wildguard} detect harmful prompts independently of the target model. Adversarial prompt shields \cite{kim2024robust} aim for robustness against jailbreaks. HiddenDetect \cite{jiang2025hiddendetect} monitors hidden states to detect attacks but requires per-model probe fitting. These approaches are either model-agnostic (ignoring fine-tuning effects) or per-model (not transferable).

\paragraph{Hypernetworks.}
Hypernetworks \cite{ha2017hypernetworks, chauhan2024brief} are neural networks that generate the weights of a target network conditioned on some context. They have been applied to continual learning \cite{vonoswald2020continual}, multi-task adaptation \cite{navon2023equivariant}, and parameter-efficient fine-tuning. Our work is, to our knowledge, the first to use hypernetworks for generating safety-specific modules conditioned on model activations.

\paragraph{Side Networks and Ladder Tuning.}
Ladder Side-Tuning \cite{sung2022lst} introduces a small side network connected to a frozen backbone via lateral (ladder) connections, enabling parameter-efficient transfer. We repurpose this architecture as a safety classifier, where the side network processes backbone representations to produce a safety score. Unlike standard LST which trains the side network per task, our hypernetwork generates side network weights for unseen models without training.


% ===========================================================================
% 3. METHODOLOGY
% ===========================================================================
\section{Methodology}
\label{sec:method}

\subsection{Problem Formulation}

Let $\mathcal{M}_0$ denote a safety-aligned base LLM and $\mathcal{M}_d$ a variant fine-tuned on domain $d$. Fine-tuning degrades safety: $\mathcal{M}_d$ may comply with harmful prompts that $\mathcal{M}_0$ would refuse. Our goal is to learn a generator $G$ such that, given any fine-tuned model $\mathcal{M}_{d^*}$ (potentially from an unseen domain $d^*$), we can produce a safety classifier $f_\theta$ that scores each input prompt $x$ as safe or harmful in the context of $\mathcal{M}_{d^*}$:
\begin{equation}
    \theta = G\big(\mathbf{A}(\mathcal{M}_{d^*})\big), \quad \lambda = f_\theta(x \mid \mathcal{M}_{d^*})
\end{equation}
where $\mathbf{A}(\mathcal{M}_{d^*})$ denotes activations extracted from $\mathcal{M}_{d^*}$ on a small calibration set and $\lambda \in [0, 1]$ is the predicted safety score ($\lambda \approx 0$: safe, $\lambda \approx 1$: harmful).

\subsection{Safety Head Architecture}
\label{sec:side_network}

We build on the Ladder Side-Tuning (LST) architecture \cite{sung2022lst}, which attaches a small transformer side network to a frozen backbone via lateral connections. The original LST was designed for parameter-efficient task adaptation, where the side network always contributes to the output. We repurpose and extend LST for safety classification with three key modifications: (1) we add a \textbf{lambda classifier head} that produces a scalar safety score from side network hidden states, (2) we introduce \textbf{selective fusion}---at inference time, the side network output is only used when the classifier predicts a harmful prompt ($\lambda > 0.5$), otherwise the base model runs unmodified, and (3) we use \textbf{selective training}---harmful prompts receive both language modeling loss and lambda classification loss, while safe prompts receive only lambda loss, so the side network focuses on safety rather than task performance.

\paragraph{Side Transformer.} The side network consists of $K=12$ transformer layers with hidden dimension $h_s = h / r$ where $h$ is the backbone hidden dimension and $r=4$ is the reduction factor. For Qwen2-7B ($h=3584$), this yields $h_s = 896$ with 7 attention heads; for LLaMA-3-8B ($h=4096$), $h_s = 1024$ with 8 heads. The total side network contains approximately 237M parameters for Qwen2 and a comparable count for LLaMA-3.

\paragraph{Ladder Connections.} At each side layer $k$, the backbone hidden state $\mathbf{h}^{(l)}_\text{back}$ from corresponding backbone layer $l$ is downsampled and fused with the side hidden state via a gated connection:
\begin{equation}
    \mathbf{h}^{(k)}_\text{in} = \mu_k \cdot W_\text{down} \mathbf{h}^{(l)}_\text{back} + (1 - \mu_k) \cdot \mathbf{h}^{(k-1)}_\text{side}
\end{equation}
where $\mu_k = \sigma(\alpha_k / \tau)$ is a learnable gate and $W_\text{down} \in \mathbb{R}^{h_s \times h}$ is the downsample projection.

\paragraph{Lambda Classifier.} After the final side layer, hidden states are mean-pooled across the sequence and passed through an MLP to produce the scalar safety score:
\begin{equation}
    \lambda = \sigma\!\left(\text{MLP}\!\left(\frac{1}{T}\sum_{t=1}^{T} \mathbf{h}^{(K)}_{t}\right) / \tau \right)
\end{equation}
where $\sigma$ is the sigmoid function and $\tau$ is a temperature scaling parameter that controls the sharpness of the safety decision boundary. We set $\tau = 2.0$ throughout, which provides a smooth yet decisive separation. At inference time, $\lambda > 0.5$ flags the prompt as harmful; for safe prompts ($\lambda \leq 0.5$), the side network is bypassed entirely and the base model output is used directly, preserving task performance.

\subsection{Activation Extraction}
\label{sec:activations}

For each fine-tuned model $\mathcal{M}_d$, we extract activations by running $N_\text{cal} = 50$ calibration prompts (domain-specific) through the model and collecting the hidden state at each transformer layer after the final token position:
\begin{equation}
    \mathbf{a}_l^{(d)} = \frac{1}{N_\text{cal}} \sum_{i=1}^{N_\text{cal}} \mathbf{h}_l^{(d)}(x_i) \in \mathbb{R}^{h}, \quad l = 1, \ldots, L
\end{equation}
These layer-wise activation vectors serve as a compact behavioral fingerprint of the fine-tuned model: they encode how fine-tuning has shifted representations relative to the base model, without requiring explicit comparison to $\mathcal{M}_0$.

\subsection{Hypernetwork Architecture}
\label{sec:hypernetwork}

The hypernetwork $G$ maps the activation fingerprint $\{\mathbf{a}_l\}_{l=1}^L$ to the side network parameters $\theta$. Directly predicting all $\sim$237M parameters from $L$ activation vectors is severely ill-conditioned. We decompose this problem through two complementary strategies.

\subsubsection{Layer-wise Factorized Generation}

Rather than mapping all activations to all side network weights simultaneously, we predict each side network layer's parameters independently from corresponding backbone activations. For side layer $k$ connected to backbone layer $l(k)$:
\begin{equation}
    \theta_k = H_k(\mathbf{a}_{l(k)})
\end{equation}
where $H_k$ is a per-layer generation module that shares most parameters across layers but includes layer-specific adapters. This reduces the effective output dimensionality from $|\theta|$ to $|\theta_k|$ per prediction.

\subsubsection{Direction-Magnitude Decomposition}

We decompose each activation vector into a unit direction and a scalar magnitude:
\begin{equation}
    \hat{\mathbf{d}}_l = \frac{\mathbf{a}_l}{\|\mathbf{a}_l\|}, \quad m_l = \log \|\mathbf{a}_l\|
\end{equation}
The direction captures \textit{what} the fine-tuning changed in representation space, while the magnitude captures \textit{how much}. These are processed by separate encoders:
\begin{align}
    \mathbf{z}_\text{dir} &= \text{Attn}\big(W_\text{proj}\, \hat{\mathbf{d}}_l\big), \quad W_\text{proj} \in \mathbb{R}^{p \times h} \\
    \mathbf{z}_\text{mag} &= \text{MLP}_\text{mag}(m_l)
\end{align}
where $p = 256$ is the direction projection dimension. The fused representation $[\mathbf{z}_\text{dir}; \mathbf{z}_\text{mag}]$ is then passed to a shared weight generation MLP that outputs the side network weights for the corresponding layer.

\subsubsection{Lambda Classifier Generation}

The lambda classifier head is generated via a dedicated branch of the hypernetwork. The fused activation embedding is first processed by a cross-layer attention module that aggregates information across all backbone layers, producing a global representation of the model's behavioral shift. This global embedding is then mapped through an MLP to produce the classifier weights. To improve efficiency, we use a low-rank factorization for the output layer:
\begin{equation}
    W_\text{cls} = B \cdot A
\end{equation}
where $A \in \mathbb{R}^{r \times h_s}$ and $B \in \mathbb{R}^{1 \times r}$ with rank $r = 16$, reducing the number of generated parameters while preserving expressiveness.

\subsection{Ground-Truth Side Network Collection}
\label{sec:side_training}

The hypernetwork requires supervised (model, side network) pairs as training data. We collect these by independently training a side network for each of the 22 fine-tuning domains. This phase is a one-time cost that produces the ground-truth targets $\{(\mathcal{M}_d, \theta_d^*)\}_{d=1}^{K}$.

\paragraph{Procedure.} For each domain $d$: (1) we fine-tune the base model $\mathcal{M}_0$ on domain data using LoRA and merge the adapter, yielding $\mathcal{M}_d$; (2) we freeze $\mathcal{M}_d$ entirely and attach a randomly initialized side network; (3) we train the side network on a mixture of domain-specific safe data (from the fine-tuning dataset) and harmful data from BeaverTails \cite{ji2024beavertails} (category-balanced, $\sim$7k samples). The side network learns to detect harmful prompts while leaving safe domain prompts unaffected.

Because all 22 side networks are initialized from the same random seed and trained with the same procedure, their weight spaces are naturally aligned---there is no permutation ambiguity---making direct weight-space comparison (and hypernetwork supervision) well-defined.

\paragraph{Selective Loss.} We apply a conditional training objective based on prompt safety:
\begin{equation}
    \mathcal{L}_\text{side} = \begin{cases}
        \mathcal{L}_\text{SFT} + \gamma \cdot \mathcal{L}_\text{cls} & \text{if } x \text{ is harmful} \\
        \gamma \cdot \mathcal{L}_\text{cls} & \text{if } x \text{ is safe}
    \end{cases}
\end{equation}
where $\mathcal{L}_\text{SFT}$ is the standard causal language modeling loss and $\mathcal{L}_\text{cls}$ is binary cross-entropy on the lambda prediction. Safe prompts receive only the classification loss because the backbone already handles them correctly---the side network only needs to learn \textit{when} to intervene, not how to generate safe responses. Harmful prompts receive both losses, teaching the side network to simultaneously detect harm and produce refusal responses.

\paragraph{Setup.} The backbone is fully frozen; only the side network ($\sim$3\% of total parameters), lambda classifier, and projection layers are trained. We use learning rate $10^{-3}$ for the side network and $10^{-4}$ for the lambda classifier, with cosine scheduling over 5--6 epochs. At inference time, the lambda score determines routing: $\lambda > 0.5$ triggers side network fusion, while $\lambda \leq 0.5$ bypasses it entirely. All $K$ domain-specific side networks trained in this phase serve as ground-truth targets for the hypernetwork.

\subsection{Hypernetwork Training}
\label{sec:training}

The hypernetwork is trained on $K$ fine-tuning domains with ground-truth side networks. We hold out $M$ domains for evaluation. Because all ground-truth side networks are trained from the same initialization, their weight spaces are aligned---there is no permutation ambiguity---making direct weight-space supervision well-defined.

\paragraph{Losses.} The training objective combines weight-space and function-space supervision:
\begin{equation}
    \mathcal{L} = \mathcal{L}_\text{recon} + \alpha \cdot \mathcal{L}_\text{lambda} + \beta \cdot \mathcal{L}_\text{func}
\end{equation}
Each loss targets a different aspect of the generation problem:
\begin{itemize}
    \item $\mathcal{L}_\text{recon}$: MSE between generated and ground-truth side network weights. This provides a dense learning signal across all $\sim$237M output dimensions, anchoring the generated weights to a known-good solution.
    \item $\mathcal{L}_\text{lambda}$: Binary cross-entropy on $\lambda$ predictions using the \textit{same} $N_\text{cal}$ calibration prompts used for activation extraction. This directly supervises the lambda classifier head on the data distribution available at inference time.
    \item $\mathcal{L}_\text{func}$: Behavioral loss on a \textit{separate} held-out set of harmful and safe probe prompts (not used for activation extraction). This ensures the generated side network generalizes beyond calibration data to unseen prompts.
\end{itemize}
The three losses are complementary: $\mathcal{L}_\text{recon}$ regularizes the high-dimensional weight space, while $\mathcal{L}_\text{lambda}$ and $\mathcal{L}_\text{func}$ ensure behavioral correctness on seen and unseen prompts respectively. We ablate each component in Section~\ref{sec:experiments}. Notably, removing $\mathcal{L}_\text{func}$ (safe probes) causes complete collapse of the domain classifier, demonstrating that behavioral supervision is critical for generalization. We set $\alpha = 10$ and $\beta = 5$ to approximately equalize gradient magnitudes. Training runs for 100 epochs with learning rate $5 \times 10^{-5}$ and hidden dimension 512.

\subsection{Inference}
\label{sec:inference}

At deployment, given a new fine-tuned model $\mathcal{M}_{d^*}$, we extract layer-wise activations from $N_\text{cal}$ domain-specific calibration prompts and pass them through the trained hypernetwork in a single forward pass to obtain the safety head weights. The generated safety head is then loaded alongside $\mathcal{M}_{d^*}$: for each input prompt, it computes $\lambda$ and flags prompts with $\lambda > 0.5$ as harmful. The entire process requires no safety data, no gradient computation, and no model modification, completing in under one minute on a single GPU.


% ===========================================================================
% 4. EXPERIMENTS
% ===========================================================================
\section{Experiments}
\label{sec:experiments}

\subsection{Setup}

\paragraph{Models and Domains.} We evaluate on Qwen2-7B-Instruct \cite{yang2024qwen2} and LLaMA-3-8B-Instruct \cite{dubey2024llama}. We construct 22 fine-tuning domains spanning mathematical reasoning (GSM8K, MATH), commonsense (CommonsenseQA, HellaSwag, PIQA, SocialIQA), factual QA (NQ-Open, BoolQ, SQuAD, ARC, PubMedQA), knowledge (MMLU, Platypus), classification (AG News), code (CodeFeedback), and instruction-following (Dolly, OpenOrca, Alpaca, FLAN, OpenHermes, NoRobots, WizardLM). Each domain is fine-tuned with LoRA (rank 16, $\alpha = 32$) and merged into the base model. Full domain details are in Appendix~\ref{app:domains}.

\paragraph{Safety Evaluation.} We use three harmful benchmarks: BeaverTails-Evaluation \cite{ji2024beavertails} (700 prompts, 14 categories), AdvBench \cite{zou2023universal} (520 prompts), and HEx-PHI \cite{qi2024hexphi} (300 prompts, 11 categories). Only BeaverTails is used during training; the other two test zero-shot transfer.

\paragraph{Metrics.} We report \textbf{Task Accuracy} (domain test set performance) and \textbf{Safety Score} (percentage of harmful prompts correctly refused, evaluated on all three benchmarks: BeaverTails, AdvBench, and HEx-PHI; higher is safer). Harmful outputs are judged by the beaver-dam-7b moderation model for BeaverTails, and by keyword-based refusal detection for AdvBench and HEx-PHI.

\paragraph{Evaluation Protocol.} We use three holdout splits, each removing four domains from different task categories. The hypernetwork trains on the remaining 18 domains and is evaluated on held-out domains it has never seen:

\begin{table}[h]
\centering
\small
\begin{tabular}{@{}cl@{}}
\toprule
\textbf{Split} & \textbf{Held-out Domains (Category)} \\
\midrule
A & BoolQ (QA), AG News (Classif.), SocialIQA (Commonsense), Comp.\ Math (Math) \\
B & CommonsenseQA (Commonsense), PIQA (Physical), GSM8K (Math), Alpaca (Instr.) \\
C & NQ-Open (QA), ARC (Science), HellaSwag (Commonsense), OpenHermes (Instr.) \\
\bottomrule
\end{tabular}
\caption{Holdout splits. Each removes domains from multiple task categories to test cross-category generalization.}
\label{tab:splits}
\end{table}

\paragraph{Baselines.} We compare against (1) the \textbf{base model} (safety-aligned, no fine-tuning), (2) the \textbf{LoRA-only} model (fine-tuned, no safety mechanism), and report both task accuracy and safety metrics.

\subsection{Main Results}

Table~\ref{tab:main_results} presents our main results across three holdout splits. Across all splits and both model families, \method{} reduces harmful output rates from 6--43\% (LoRA-only) to near 0\% across all three safety benchmarks, while preserving task accuracy within 1 percentage point of the fine-tuned model. Fine-tuning introduces substantial safety degradation (e.g., PIQA on Qwen2: 43\% harmful rate on BeaverTails), and our generated safety head eliminates this degradation without sacrificing the task performance gains. Notably, the safety recovery transfers zero-shot to AdvBench and HEx-PHI, which are never seen during hypernetwork training.

\begin{table*}[t]
\centering
\resizebox{\textwidth}{!}{%
\begin{tabular}{@{}cl ccc ccc ccc ccc@{}}
\toprule
& & \multicolumn{6}{c}{\textbf{LLaMA-3-8B}} & \multicolumn{6}{c}{\textbf{Qwen2-7B}} \\
\cmidrule(lr){3-8} \cmidrule(lr){9-14}
& & \multicolumn{3}{c}{Task Acc (\%)} & \multicolumn{3}{c}{Safety (\%) $\uparrow$} & \multicolumn{3}{c}{Task Acc (\%)} & \multicolumn{3}{c}{Safety (\%) $\uparrow$} \\
\cmidrule(lr){3-5} \cmidrule(lr){6-8} \cmidrule(lr){9-11} \cmidrule(lr){12-14}
\textbf{Split} & \textbf{Domain} & Base & LoRA & Ours & BT & Adv & HEx & Base & LoRA & Ours & BT & Adv & HEx \\
\midrule
\multirow{4}{*}{A}
& BoolQ       & 65.8 & 92.2 & 92.1\worse{$_{-0.1}$} & 99.9 & 100 & 100  & 81.9 & 92.9 & 92.8\worse{$_{-0.1}$} & 99.4 & 100 & 98.0 \\
& AG News     & 77.6 & 90.6 & 90.6\better{$_{\pm0}$} & 99.7 & 100 & 99.7 & 72.1 & 98.7 & 98.7\better{$_{\pm0}$} & 99.7 & 100 & 98.0 \\
& SocialIQA   & 77.8 & 97.0 & 96.8\worse{$_{-0.2}$} & 99.4 & 100 & 99.3 & 70.2 & 92.5 & 92.3\worse{$_{-0.2}$} & 99.1 & 100 & 97.3 \\
& Comp.\ Math & 32.4 & 34.1 & 34.1\better{$_{\pm0}$} & 99.6 & 100 & 99.7 & 42.2 & 74.4 & 74.4\better{$_{\pm0}$} & 99.3 & 100 & 98.7 \\
\midrule
\multirow{4}{*}{B}
& CSenseQA    & 70.8 & 97.1 & 96.8\worse{$_{-0.3}$} & 99.4 & 100 & 99.7 & 72.7 & 94.0 & 93.8\worse{$_{-0.2}$} & 99.1 & 100 & 97.3 \\
& PIQA        & 60.4 & 97.5 & 97.2\worse{$_{-0.3}$} & 99.7 & 100 & 100  & 56.7 & 99.1 & 99.0\worse{$_{-0.1}$} & 99.4 & 100 & 98.3 \\
& GSM8K       & 76.9 & 67.3 & 67.3\better{$_{\pm0}$} & 99.0 & 100 & 99.3 & 79.0 & 76.2 & 76.2\better{$_{\pm0}$} & 98.9 & 100 & 97.7 \\
& Alpaca      & 24.8 & 39.0 & 38.8\worse{$_{-0.2}$} & 99.3 & 100 & 99.3 & 29.5 & 39.6 & 39.5\worse{$_{-0.1}$} & 99.1 & 100 & 97.7 \\
\midrule
\multirow{4}{*}{C}
& NQ-Open     & 35.6 & 77.5 & 77.0\worse{$_{-0.5}$} & 98.7 & 100 & 99.0 & 30.3 & 54.8 & 54.6\worse{$_{-0.2}$} & 99.1 & 100 & 97.3 \\
& ARC         & 87.7 & 86.5 & 86.5\better{$_{\pm0}$} & 99.3 & 100 & 99.7 & 89.2 & 91.3 & 91.3\better{$_{\pm0}$} & 99.0 & 100 & 98.0 \\
& HellaSwag   & 50.9 & 92.2 & 91.5\worse{$_{-0.7}$} & 99.7 & 100 & 100  & 67.0 & 97.8 & 97.4\worse{$_{-0.4}$} & 98.6 & 100 & 97.7 \\
& OpenHermes  & 24.8 & 54.8 & 54.5\worse{$_{-0.3}$} & 99.4 & 100 & 99.7 & 35.9 & 57.2 & 57.0\worse{$_{-0.2}$} & 99.0 & 100 & 97.3 \\
\midrule
\multicolumn{2}{@{}l}{\textbf{Average}} & --- & --- & \worse{$_{-0.2}$} & 99.4 & 100 & 99.6 & --- & --- & \worse{$_{-0.1}$} & 99.1 & 100 & 97.8 \\
\bottomrule
\end{tabular}%
}
\caption{Main results across three holdout splits on unseen domains. \textbf{Task Acc}: domain-specific accuracy. \textbf{Safety}: harmful prompt detection rate on BeaverTails (BT, 700 prompts), AdvBench (Adv, 520 prompts), and HEx-PHI (HEx, 300 prompts). Only BT is used during training; Adv and HEx test zero-shot transfer. Subscripts show accuracy change vs.\ LoRA ({\color{green!60!black}green}: preserved, {\color{red!70!black}red}: minor drop). Each split holds out 4 of 22 domains; the hypernetwork trains on the remaining 18.}
\label{tab:main_results}
\end{table*}

\subsection{Ablation Studies}

We conduct ablation studies on LLaMA-3-8B to validate our design choices. We report two complementary metrics: \textbf{Task Preservation} (percentage of safe domain prompts correctly passed through without side network intervention, higher is better) and \textbf{Safety Recovery} (percentage of harmful prompts correctly detected, higher is better). Results are averaged across four holdout domains (BoolQ, CommonsenseQA, PIQA, SocialIQA).

\paragraph{Loss Components.} Table~\ref{tab:ablation_loss} ablates each loss term. Removing $\mathcal{L}_\text{func}^\text{safe}$ (safe probe supervision) is catastrophic: the classifier marks \textit{all} inputs as harmful, achieving perfect safety but zero task preservation. Removing all functional supervision ($-\mathcal{L}_\text{func}$) degrades task preservation by 20+ points, showing that behavioral supervision is essential for the classifier to distinguish safe from harmful inputs. Domain-specific calibration prompts are critical for safety: without them, task preservation improves (the classifier defaults to ``safe'') but safety recovery drops sharply.

\begin{table}[t]
\centering
\small
\begin{tabular}{@{}lcc@{}}
\toprule
\textbf{Variant} & \textbf{Task Pres.\ (\%) $\uparrow$} & \textbf{Safety Rec.\ (\%) $\uparrow$} \\
\midrule
Full system & 83.7 & 99.5 \\
$-\,\mathcal{L}_\text{func}^\text{safe}$ & 6.5 & 100.0 \\
$-\,\mathcal{L}_\text{func}$ (recon only) & 61.4 & 99.8 \\
$-\,$Domain calibration & 98.1 & 89.3 \\
$-\,$Lambda LoRA gen. & 57.7 & 100.0 \\
\bottomrule
\end{tabular}
\caption{Loss and component ablation on LLaMA-3. Task Preservation: safe prompts correctly passed through. Safety Recovery: harmful prompts correctly detected. Averaged over four holdout domains.}
\label{tab:ablation_loss}
\end{table}

\paragraph{Training Data Composition.} Table~\ref{tab:ablation_data} examines which categories of training domains are most important. Removing instruction-following domains (6 domains) causes a dramatic shift: the classifier defaults to ``safe,'' achieving near-perfect task preservation but losing 5+ points of safety recovery. This reveals that instruction-following data provides the diversity needed to learn the safe-vs-harmful boundary. Removing QA or math/code domains has minimal impact on safety, and removing knowledge domains (MMLU, Platypus, AG~News) actually \textit{improves} task preservation, suggesting these domains add noise to the activation fingerprint rather than useful signal.

\begin{table}[t]
\centering
\small
\begin{tabular}{@{}lccc@{}}
\toprule
\textbf{Excluded Category} & \textbf{\#Excl.} & \textbf{Task Pres.} & \textbf{Safety Rec.} \\
\midrule
None (full, 18 domains) & 0 & 83.7 & 99.5 \\
$-\,$QA domains & 3 & 80.3 & 100.0 \\
$-\,$Math \& code & 3 & 84.8 & 99.7 \\
$-\,$Knowledge \& classif. & 3 & 87.8 & 99.5 \\
$-\,$Instruction-following & 6 & 99.6 & 91.2 \\
\bottomrule
\end{tabular}
\caption{Effect of excluding entire training data categories. Instruction-following domains are critical for safety discrimination. Removing knowledge domains slightly improves task preservation without harming safety.}
\label{tab:ablation_data}
\end{table}

\paragraph{Domain Exclusion Robustness.} To test robustness to reduced training diversity, we progressively exclude 1--2 individual domains. Excluding 1 domain has negligible impact ($<$0.5\% change in both metrics); excluding 2 domains from different categories shows $<$1\% variation, confirming the hypernetwork's robustness to minor training set variations.

\subsection{Analysis}

\paragraph{Cross-Benchmark Transfer.} A key result in Table~\ref{tab:main_results} is the consistent safety performance across all three benchmarks. The hypernetwork is trained with BeaverTails as the sole harmful data source, yet generated safety heads achieve 100\% detection on AdvBench and 97--100\% on HEx-PHI across all holdout domains and both models. This confirms that the classifier captures general harmful content patterns rather than BeaverTails-specific features. The consistency across 12 unseen holdout domains further demonstrates that the safety heads generalize across both model architectures and task categories.

\paragraph{Lambda Score Distributions.} Figure~\ref{fig:lambda_dist} shows $\lambda$ score distributions for safe vs.\ harmful prompts on representative unseen domains. Safe domain prompts cluster near $\lambda \approx 0$ while harmful prompts cluster near $\lambda \approx 1$, with clear separation at the $\lambda = 0.5$ threshold. This bimodal distribution confirms that the hypernetwork-generated classifier learns a meaningful safe--harmful boundary rather than relying on a fragile threshold.


% ===========================================================================
% 5. CONCLUSIONS
% ===========================================================================
\section{Conclusions}
\label{sec:conclusions}

We introduced \method{}, a framework for model-conditioned safety assessment that generates lightweight safety classifiers for fine-tuned LLMs at inference time. By training a hypernetwork to map model activations to safety head weights, we amortize the cost of per-model safety restoration into a single training phase. Generated safety heads recover harmful prompt detection on unseen fine-tuning domains while preserving task performance ($<$1\% accuracy drop), validated across two model architectures (Qwen2-7B, LLaMA-3-8B), three holdout splits covering 12 diverse domains, and three safety benchmarks (BeaverTails, AdvBench, HEx-PHI). Notably, the safety classifier trained only on BeaverTails transfers to AdvBench (100\% detection) and HEx-PHI (97--100\%), demonstrating that the hypernetwork learns generalizable safety patterns rather than benchmark-specific features. Our ablation studies reveal that instruction-following training data and safe-probe behavioral supervision are the two most critical ingredients for effective safety head generation. Our approach offers a scalable path toward safety-as-a-service for the growing ecosystem of fine-tuned language models.

\paragraph{Future Work.} Extending to generation-level safety monitoring (beyond prompt classification), scaling to larger models (70B+), exploring cross-architecture transfer, and evaluating robustness against adversarial jailbreak attacks are promising directions.


\section*{Limitations}

Our approach has several limitations. First, the hypernetwork training requires ground-truth side networks for multiple domains, which involves a one-time but non-trivial setup cost. Second, the lambda classifier operates at the prompt level and does not monitor generated outputs, so harmful continuations of safe-seeming prompts may not be caught. Third, we evaluate on 7B--8B parameter models; scaling behavior to significantly larger models remains untested. Fourth, the side network adds approximately 3\% additional parameters and a corresponding increase in inference latency. Fifth, while we evaluate across three harmful benchmarks (BeaverTails, AdvBench, HEx-PHI), evaluation on adversarial jailbreak attacks and multi-turn conversations would further validate robustness. Finally, we do not evaluate in multilingual settings.


% ===========================================================================
% REFERENCES
% ===========================================================================
\bibliography{custom}

% ===========================================================================
% APPENDIX
% ===========================================================================
\appendix

\section{Domain Dataset Details}
\label{app:domains}

\begin{table}[h]
\centering
\small
\begin{tabular}{@{}llrl@{}}
\toprule
\textbf{Domain} & \textbf{Category} & \textbf{Train} & \textbf{Source} \\
\midrule
BoolQ & QA & 9.4k & \cite{clark2019boolq} \\
NQ-Open & QA & 10k & \cite{lee2019latent} \\
SQuAD & QA & 26k & \cite{rajpurkar2016squad} \\
PubMedQA & QA & 0.5k & \cite{jin2019pubmedqa} \\
ARC & Science QA & 2.4k & \cite{clark2018think} \\
CommonsenseQA & Commonsense & 9.7k & \cite{talmor2019commonsenseqa} \\
PIQA & Physical & 16.1k & \cite{bisk2020piqa} \\
HellaSwag & Commonsense & 10k & \cite{zellers2019hellaswag} \\
SocialIQA & Social & 10k & \cite{sap2019socialiqa} \\
GSM8K & Math & 7.5k & \cite{cobbe2021training} \\
Comp.\ Math & Math & 7.5k & \cite{hendrycks2021measuring} \\
CodeFeedback & Code & 10k & \cite{zheng2024opencodeinterpreter} \\
AG News & Classification & 24k & \cite{zhang2015character} \\
MMLU & Knowledge & 14k & \cite{hendrycks2021test} \\
Platypus & Knowledge & 5k & \cite{lee2023platypus} \\
Dolly & Instruction & 15k & \cite{conover2023free} \\
OpenOrca & Instruction & 10k & \cite{OpenOrca} \\
Alpaca & Instruction & 52k & \cite{taori2023stanford} \\
FLAN & Instruction & 10k & \cite{wei2022finetuned} \\
OpenHermes & Instruction & 10k & \cite{OpenHermes} \\
NoRobots & Instruction & 10k & \cite{norobots} \\
WizardLM & Instruction & 10k & \cite{xu2024wizardlm} \\
\bottomrule
\end{tabular}
\caption{The 22 fine-tuning domains used across both Qwen2-7B and LLaMA-3-8B. All models are fine-tuned with LoRA (rank 16, $\alpha = 32$) and merged.}
\label{tab:domains}
\end{table}

\section{Implementation Details}
\label{app:implementation}

\paragraph{Side Network.} For Qwen2-7B: hidden dimension 896, 7 attention heads, intermediate 4736, 12 layers ($\sim$237M parameters, 2.88\% of backbone). For LLaMA-3-8B: hidden dimension 1024, 8 attention heads, intermediate 3584, 12 layers.

\paragraph{LoRA Fine-Tuning.} Rank 16, $\alpha = 32$, applied to all linear layers. Trained for 3 epochs with learning rate $2 \times 10^{-4}$.

\paragraph{Hypernetwork Training.} Hidden dimension 512, lambda LoRA rank 16, learning rate $5 \times 10^{-5}$, 100 epochs, $N_\text{cal} = 50$ calibration prompts per domain, direction projection dimension 256, magnitude hidden dimension 128. Loss weights: $\alpha = 10$ (lambda), $\beta = 5$ (functional).

\section{Per-Domain Results}
\label{app:per_domain}

\remark{Full tables for all 4 splits $\times$ all domains $\times$ both architectures.}

\section{Additional Ablations}
\label{app:ablations}

\remark{Extended ablation results, including sensitivity to hidden dimension, rank, number of training domains.}

\end{document}
